\documentclass[10pt]{article}
\usepackage[a4paper,margin=2.4cm]{geometry}
\usepackage{array}
\usepackage{booktabs}
\usepackage{amsmath}
\usepackage{amssymb}
\usepackage[hidelinks]{hyperref}
\usepackage{xcolor}
\usepackage{float}
\setlength{\parskip}{2pt}
\newcommand{\code}[1]{\texttt{#1}}
\newcolumntype{L}[1]{>{\raggedright\arraybackslash}p{#1}}

\title{Family 0.9.tf: land and coast through a published embedding\\
\large a dataset design note --- family 1.0.tf with its raw imagery replaced by Google's AlphaEarth annual 10\,m embeddings}
\author{blauewelt/earth --- design note}
\date{16 September 2026}

\begin{document}
\maketitle

\begin{abstract}
Family 0.9.tf (``tf'' for \emph{terrestrial, fine}; a \emph{family} is a
numbered set of datasets built to one storage design and read by the
programme's forecasting model) is family 1.0.tf --- the land-and-coast
family of raw fine-scale observations --- with \textbf{one substitution}.
Family 1.0.tf's tier-T (image-tile) scene catalogue (its ten image-catalogue stores,
\code{cat\_landsat} to \code{cat\_submetre}: the raw satellite imagery and
the sub-metre orthophotos, listed for a small learned encoder --- the
\emph{local codec} --- to turn into tokens) is replaced by four stores of
the AlphaEarth Satellite Embedding, Google DeepMind's published summary of
that imagery: \code{aef1k}, \code{aef\_dots}, \code{aef100r} and
\code{aef10}. Everything else --- the point networks, the LiDAR footprints,
the daily 0.05$^\circ$ fields, the read-out targets and the statics --- is
\textbf{inherited from 1.0.tf by reference}, exactly as family 10.2
inherits family 10.1's stores: the 0.9.tf registry lists them with their
1.0.tf paths and sha256 checksums, and not one byte is rebuilt. So
0.9.tf's own build is the four AlphaEarth stores, and the comparison the two
families exist for --- raw imagery through a learned local codec versus a
published embedding, everything else held fixed --- is exact.

Instead of learning a local codec over petabytes of raw optical, radar and
LiDAR (laser-ranging) scenes, the family ingests the embedding Google already
computed from the same optical and radar imagery, trained against the same
LiDAR, L-band radar (long-wavelength, about 24\,cm) and climate sources the
catalogue lists --- the AlphaEarth Foundations Satellite Embedding, version 1:
one 64-dimensional unit-length vector per 10\,m pixel per year, 2017--2025,
for every land surface and shallow water, published under CC BY 4.0 (free to use with attribution), 64\,TB
a year and 577\,TB in all, readable anonymously from a complete mirror. The
version number below 1.0 says what it is: a published summary standing in
for the raw record, the fastest route to ten-metre land information the
programme can take, bought at the price of an annual cadence (73 times
coarser than the five-day rule), a closed model that cannot be re-run on new
sensors or sub-annual periods, and a record that begins in 2017. The note
says what 0.9.tf inherits, what the embedding is and what it was trained on, how it enters family
10's token (the unit the model reads: a value, its offset in space and time, and its footprint) without leaking a year's data into its own anchors, three stores
at three resolutions that make it usable at 0.2\,TB in the first phase
rather than 577 (a 1\,km
pooled field with a coherence channel, raw 10\,m samples at fixed sub-cell
positions, and 100\,m over regions of interest), the 10\,m original
referenced by byte range rather than copied, the arithmetic, the one-pass
build near the data, the falsifiers that decide it against family 1.0.tf's
own codec, and the decisions left to Chris.
\end{abstract}

\section{Purpose --- what this family is for}

Family 1.0.tf holds the land's raw fine record as a catalogue of scenes, and
its imagery becomes usable only once a local codec (E-078 \S4) has been
trained to turn tiles into tokens --- the one compute-bound rung on the
programme's ladder, and one whose first build is coastal and small
(E-078c, the first experiment to train that codec, on about 300 coastal cells; E-078 is the programme's multi-granularity design). Google's AlphaEarth Foundations model has already done that job
at planetary scale: it reads the same optical and C-band radar imagery the
catalogue lists (Sentinel-2, Sentinel-1, Landsat 8 and 9 --- the EU's optical and radar satellites and the US optical ones) and was trained
to reproduce the rest of it --- L-band radar, spaceborne LiDAR, climate and
gravity fields, elevation, land-cover labels and geo-linked text --- as
targets, so those sources shape the embedding without being read when it
runs; and its annual output is public. Family 0.9.tf takes that output as
a token source in place of 1.0.tf's image catalogue, keeps every other 1.0.tf store unchanged (\S\ref{sec:inherits}), so that the land side of ``predict everything'' can be
tested \emph{now}, against the land targets of family 1.0.tf (forest
disturbance, loss year, biomass, burned area, water extent, crop and land
cover), and so that the programme can measure whether its own codec is
needed for land at all.

The number 0.9 is deliberate. It is below 1.0 because the source is not an
observation but a learned summary of observations; the family is a
stepping stone, and it is meant to be \emph{compared} with 1.0.tf, not
merged into it; because the two share every other store byte for byte, the comparison isolates the imagery. If AlphaEarth tokens predict the land targets as well as the
programme's own codec would, the codec is not worth its compute for land at
annual cadence; if they do not --- because the targets move faster than a
year, or because the closed model cannot see NISAR (the NASA--ISRO L-band radar satellite launched in 2025), or because 2017 is too
late --- the comparison says exactly where 1.0.tf's raw record earns its
cost.

\section{What 0.9.tf inherits from 1.0.tf --- every store except the imagery}
\label{sec:inherits}

Family 0.9.tf builds four stores and borrows every other one. It inherits
family 1.0.tf's stores \emph{by reference}, the way family 10.2 inherits
family 10.1's: the 0.9.tf registry (\S\ref{sec:stores}) lists each inherited
store with its 1.0.tf path and the sha256 checksum of every file as 1.0.tf
published it, and nothing is re-fetched, rebuilt, re-hashed or re-uploaded.
An inherited store is therefore the same bytes in both families, and a model
trained on either reads it from the same file. Source, channels, cadence,
size and licence of each inherited store are in the family 1.0.tf note
(\S4 there) and are not repeated here.

\begin{table}[H]
\centering\setlength{\tabcolsep}{3pt}\footnotesize
\begin{tabular}{@{}L{3.3cm}L{5.4cm}L{5.6cm}L{1.3cm}@{}}
\toprule
group (tier) & stores & what they are & 1.0.tf \\
\midrule
point networks (P) & \code{ghcnd}, \code{igra}, \code{tide}, \code{tide\_private}$^{*}$, \code{ndbc}, \code{gliders}, \code{flux}, \code{fire}, \code{river}, \code{river\_grdc}$^{*}$, \code{ismn}$^{*}$, \code{ecad}$^{*}$, \code{ghcnh}, \code{hfr}, \code{alticap} & weather stations, radiosonde balloons, tide gauges, buoys, gliders, flux towers, fire detections, river gauges, soil-moisture probes, coastal radar and coastal altimetry & \S4.1 \\
LiDAR footprints (P) & \code{gedi}, \code{icesat2} & laser measurements of canopy and terrain height from orbit & \S4.2 \\
daily fields at 0.05$^\circ$ (G sharded) & \code{lst05}, \code{snow05}, \code{refl05}, \code{ims1k}, \code{chirps05} & land-surface temperature, snow cover, surface reflectance, snow-and-ice charts, five-day rainfall & \S4.3 \\
read-out targets & \code{alerts}, \code{lossyear}, \code{burned500}, \code{cat\_gfm}, \code{biomass100}, \code{landcover10} & forest-clearing alerts and loss years, burned area, flood maps, biomass, land and crop cover & \S4.5 \\
statics & \code{static\_fine} & elevation and sea-floor depth, land-cover fractions, distance to coast & \S4.6 \\
\midrule
\emph{not inherited} --- replaced by \S\ref{sec:stores} here & \code{cat\_landsat}, \code{cat\_hls}, \code{cat\_s2}, \code{cat\_s1}, \code{cat\_nisar}, \code{cat\_palsar}, \code{cat\_olci}, \code{cat\_viirs}, \code{cat\_ecostress}, \code{cat\_submetre}; the tile tokens computed from them; \code{oc300\_coast} & the tier-T image catalogue --- the raw satellite imagery and the sub-metre orthophotos --- and what is derived from it & \S4.4 \\
\bottomrule
\end{tabular}
\caption{Family 1.0.tf's stores as family 0.9.tf uses them. $^{*}$ on 1.0.tf's
private track: the store's licence forbids republishing, so it stays in the
private repository and 0.9.tf's registry lists it there with
\code{distribution: private}.}
\end{table}

Two boundary cases. \code{cat\_gfm} has a catalogue's prefix because it is
one, but of flood maps --- a read-out target, not imagery --- so it is
inherited. The phase-C tile tokens and \code{oc300\_coast} (the 300\,m
gridded form of \code{cat\_olci}) are made from the imagery and go with it.
An inherited store appears in 0.9.tf's registry when 1.0.tf publishes it,
with the same checksums; 0.9.tf never waits on, re-runs or patches a 1.0.tf
build, and a fix to an inherited store is a new 1.0.tf version that both
registries then name.

\section{What the embedding is --- Google's AlphaEarth model and its published 64 numbers per 10\,m pixel per year}

Checked against the paper (Brown et al.\ 2025, arXiv --- the open preprint server --- 2507.22291, v2
2025-09-08), the catalogue entry of Earth Engine (Google's hosted geospatial analysis platform), Google's storage read-me and
a live listing of both buckets on 2026-09-16.

\paragraph{The model.} An ``embedding field'' model of about 480\,M
parameters (a 1\,B variant was trained; the released layers come from the
480\,M one). Its \emph{inputs} are image sequences from Sentinel-2,
Sentinel-1 C-band radar and Landsat 8 and 9. Its training \emph{targets},
which the model learns to reconstruct and therefore encodes without reading
them at run time, are ALOS PALSAR-2 L-band radar (Japan's long-wavelength radar satellite), GEDI LiDAR (laser canopy heights from the Space Station), ERA5-Land
climate (a land reanalysis from ECMWF, the European weather centre), GRACE gravity (satellites that weigh water through its pull), the Copernicus GLO-30 elevation model (the EU's 30\,m terrain map), NLCD
land-cover labels (the US National Land Cover Database), and Wikipedia articles linked to species
occurrences from GBIF (the Global Biodiversity Information Facility) --- more than three billion observations in all, over about
1.1\,\% of the land surface. No sub-metre imagery is used. The model separates a \emph{support period} (the
timestamps of its inputs) from a \emph{valid period} the output summarises,
and can summarise, interpolate or extrapolate to a valid period with no
observations in it; the released product uses a valid period of one calendar
year, and the exact support window of the released layers is not stated.

\paragraph{The output.} Per 10\,m pixel per year, a 64-dimensional vector
of unit length --- the vectors lie on the sphere $S^{63}$, spread by a
batch-uniformity objective --- which the producer describes as linearly
composable, with the dot product between two years at one pixel as the
documented change measure (cosine similarity, since the vectors are unit
length). Coverage is terrestrial surfaces and shallow water; years
2017--2025 (the 2025 layer was released 2026-03-09; Google commits to annual
production with a year's notice of change, contingent on Landsat and
Copernicus continuing).

\paragraph{The files.} Cloud-optimised GeoTIFFs (COGs: image files that can be read piece by piece over HTTP) of $8192 \times 8192$ pixels
with 64 bands, in local UTM projections (Universal Transverse Mercator map zones), 120 zone directories per year,
about 34,000 tiles a year, 1.9\,GB on average. Values are \textbf{int8} (one signed byte each)
with nodata (the no-data marker) $-128$ and the documented de-quantisation
\begin{equation}
x = \Big(\frac{q}{127.5}\Big)^{2}\,\mathrm{sign}(q),
\end{equation}
so the encoding is $q \approx \mathrm{round}(127.5\,\mathrm{sign}(x)\sqrt{|x|})$
--- a square-root companding that spends resolution near zero. The producer
states the 8-bit form costs negligible accuracy against float32.

\sloppy\paragraph{Where it lives, measured.} On Google Cloud Storage, \code{gs://alphaearth\_foundations/}\allowbreak
\code{satellite\_embedding/}\allowbreak\code{v1/annual/}\allowbreak\code{<year>/<zone>/}, with index files
(\code{aef\_index.parquet}, a table of every tile, 70\,MB, updated 2026-02-27); the bucket has been
provider-pays since July 2026, so a copy costs the reader nothing. A
complete mirror on Source Cooperative (a public data host; \code{tge-labs/aef}, under AWS Open Data, Amazon's sponsored open-data programme,
\code{s3://\allowbreak us-west-2.opendata.source.coop/\allowbreak tge-labs/\allowbreak aef/\allowbreak v1/\allowbreak annual/}, anonymous
S3 --- Amazon's object-storage protocol --- and HTTPS) has the same tile counts and byte totals as the origin for
every year, plus one VRT (a small GDAL virtual-raster file that says how to read a tile) per tile; its read-me warns that the COGs are
stored bottom-up, so they are read through the VRTs. Earth Engine serves the
same layers as floats in $[-1, 1]$ (bands \code{A00}--\code{A63}).\par\fussy

\begin{table}[H]
\centering\footnotesize
\begin{tabular}{@{}lrr@{}}
\toprule
year & tiles & TB \\
\midrule
2017 & 33,148 & 64.35 \\
2018 & 33,291 & 64.15 \\
2019 & 33,385 & 64.06 \\
2020 & 33,206 & 63.97 \\
2021 & 33,281 & 63.91 \\
2022 & 33,700 & 63.11 \\
2023 & 34,151 & 63.85 \\
2024 & 34,149 & 63.90 \\
2025 & 34,155 & 65.29 \\
\midrule
total & $\approx 302{,}000$ & \textbf{576.6} \\
\bottomrule
\end{tabular}
\caption{The published embedding, as listed on 2026-09-16 (origin and mirror agree year by year). About $1.4 \times 10^{12}$ land pixels a year.}
\end{table}

\paragraph{Licence.} CC BY 4.0, attribution ``The AlphaEarth Foundations
Satellite Embedding dataset is produced by Google and Google DeepMind''.
Every AlphaEarth store below carries it; \code{licence.redistribution} is
\code{attribution} and \code{distribution} is \code{public}. The two-track
rule of family 1.gf \S4 (public stores in the public repository, stores whose licence forbids republishing in a private one) applies, and none of the four AlphaEarth stores needs the
private track today; the four private stores 0.9.tf inherits stay on 1.0.tf's private track (\S\ref{sec:inherits}).

\paragraph{What is not available.} Only version 1 and only annual. Sub-annual
embeddings (quarterly, monthly, weekly, ``down to a 5-day frequency'') exist
as Google Maps Platform's Custom Satellite Embeddings, in private preview
with no public licence; an academic programme offers samples through Earth
Engine, applications closing 2026-10-15. A granted sample is usable here under the programme's terms as a
private-track store, \code{aef\_custom} (sub-annual embeddings for this project's use only; family 1.gf \S4), never
republished --- the one place this family's own build would use that track. The model and training code are
closed: no new sensor (NISAR), no new year on demand, no re-run over a
chosen period. Reported weaknesses: transfer across regions and years is
limited (an irrigation study reproduced 0.93--0.95 accuracy locally and
about 0.3 when moved to another region), short events are averaged away by
the annual window, and the model receives no coordinates, so a
location-dependent task needed far more labels than a coordinate-aware
baseline.

\section{Grid, time axis and token --- how an annual 10\,m map becomes model input without leaking the future}

The embedding is a 10\,m raster in UTM; the family's grids are latitude and
longitude. The stores below resample by \emph{pooling}, never by
interpolating a single value, and the 10\,m original is referenced, not
copied (\S\ref{sec:aef10}), so nothing is resampled at storage time except what is
declared as a pooled field with its own footprint.

\paragraph{Time.} A year-$Y$ embedding summarises 1 January to 31 December
of $Y$: $\log_2 f_t = \log_2(365/5) = +6.2$ (the time span a value summarises, as a base-2 logarithm in units of five days). It is filed under the
\textbf{first family-10 bin (five-day time step counted from 1982) that begins after 31 December $Y$} --- not the
bin containing that day, because the reader includes the anchor's own bin
(family 10.1 \S7: candidates are ``at or before the anchor's pentad''), and
a year-end bin would hand an anchor in the last days of $Y$ its own year's
summary. With that filing the one-sided rule does the rest: an anchor in
any bin of year $Y$ sees the embedding of year $Y-1$ and earlier, never
$Y$'s own. $\Delta t$ is measured from the end of the anchor's bin to
31 December $Y$, 0 to about 370 days for the most recent usable year. This
is the strict rule and the pick; the relaxed alternative --- letting an
anchor in $Y$ read $Y$'s embedding when the target lies after $Y$ --- is a
leak of up to a year of future imagery into the input and is not offered.
One caveat the rule cannot close: the producer does not state the support
window of the released layers, so if year $Y$'s embedding was computed
from observations extending past 31 December $Y$, even this filing admits
them; the registry records the assumption.

\paragraph{Footprint.} A raw 10\,m vector carries $\log_2 f_p$ (the area a value summarises, as a base-2 logarithm of 27.83\,km) $=
\log_2(0.01/27.83) = -11.4$; a pooled $1/1200^\circ$ cell (0.093\,km)
$-8.2$; a pooled $1/120^\circ$ cell (0.93\,km) $-4.9$. All are inside the
registry's $\pm 12$ clamp.

\paragraph{The token.} The 64 components are the token's value vector with
$C = 64$ channels (plus one coherence channel for pooled stores, \S\ref{sec:aef1k}), raw units
(unit-length components in $[-1, 1]$), mask all-true where the pixel is not
nodata; source embedding ``AlphaEarth v1''; offsets and footprints as
everywhere else. Nothing is z-scored (standardised to zero mean and unit variance): the components are already on a
sphere, and the loader standardises from training years if it wants to.

\section{The stores --- the four AlphaEarth datasets this family builds, and its registry}
\label{sec:stores}

\paragraph{The registry.} \code{tensors/\allowbreak family09\_tf/\allowbreak family09tf.json}
on \code{chfrank/earth-tensors} (the project's public Hugging Face dataset
repository), in the schema of family 1.gf's registry (\S4 of that note): the
same top-level keys, and a \code{path} per group so a consumer resolves every
store where it actually lives. It carries an \textbf{\code{inherits} block}
naming family 1.0.tf's registry, \code{tensors/\allowbreak family1\_tf/\allowbreak family1tf.json}:
\begin{verbatim}
"inherits": {
  "1.0.tf": {
    "root":     "tensors/family1_tf",
    "registry": "tensors/family1_tf/family1tf.json",
    "groups":   ["ghcnd", "igra", "tide", ..., "static_fine"],
    "note":     "built by family 1.0.tf and unchanged: not one
                 byte rebuilt, re-hashed or re-uploaded for 0.9.tf"
  }
}
\end{verbatim}
Each inherited group is listed with its 1.0.tf \code{path} and the
\code{sha256} of every file exactly as 1.0.tf's registry publishes them ---
the four private ones with their private-repository path and
\code{distribution: private} --- so a request under
\code{tensors/\allowbreak family09\_tf/\allowbreak ghcnd/} correctly fails: that store lives under
\code{tensors/\allowbreak family1\_tf/\allowbreak ghcnd/}. The four groups below are the only ones
whose \code{path} is under \code{tensors/\allowbreak family09\_tf/}. Family 10.2's
registry names this file in its \code{siblings} list (family 1.gf \S4).

Three stores at three resolutions, and one reference. Bytes assume int8
components (the producer's own precision) plus one uint8 coherence channel,
so a pooled pixel costs 65 bytes and a raw 10\,m vector 64.

\subsection{\code{aef1k} --- the embedding averaged over each square kilometre, with a number saying how uniform the square is (tier G sharded)}\label{sec:aef1k}

On a $1/120^\circ$ grid ($\approx 0.93$\,km, $43200 \times 21600$), each
cell's value is the \textbf{mean of the unit vectors} of its $\approx 100
\times 100$ valid 10\,m pixels, \textbf{renormalised} to unit length, stored
as int8 under the producer's companding; and a 65th channel, the
\textbf{coherence}: the length of the un-normalised mean, in $[0, 1]$, 1
where the square kilometre is one thing and near 0 where it is many. The
coherence is what the mean throws away, kept as a number, so that a 1\,km
token says both ``what'' and ``how uniformly''. One frame per year
($F = 1$), filed under the bin of 31 December. Land is $\approx 29$\,\% of
the grid, $\approx 270$\,M valid cells: $270\,\mathrm{M} \times 65 \times
1.2 \approx 21$\,GB per year, \textbf{$\approx 0.2$\,TB for 2017--2025}
(0.25\,TB if the high-latitude land the 29\,\% undercounts is included).

\subsection{\code{aef\_dots} --- sixteen unaveraged 10\,m embedding samples per quarter-degree cell per year (tier P, point rows)}

Pooling loses the fine structure the embedding was made for. \code{aef\_dots}
keeps sixteen \emph{unpooled} 10\,m vectors per $0.25^\circ$ land cell per
year at fixed sub-cell positions --- the cone's own sunflower offsets (the fixed spiral of sample points around the predicted place that the model's input cone uses) scaled
into the cell, the same sixteen positions every year so that a position's
time series is a time series --- as tier-P rows: \code{lat}, \code{lon} of
the pixel centre, \code{time\_s} = 31 December $Y$, \code{platform} = the
hashed tile identifier, \code{qc} = 1, \code{values} the 64 components as
float16 (two-byte numbers) ($C = 64$; int8 would need a schema change and is not worth it at
this size). About 300\,k land cells $\times$ 9 years $\times$ 16 =
43\,M rows $\times$ 155 bytes $\approx$ \textbf{7\,GB}. This is the store a
consumer reads to ask what 10\,m detail adds over the kilometre mean, for
7\,GB.

\subsection{\code{aef100r} --- the embedding averaged over 100\,m squares, only in chosen regions, deforestation frontiers first (tier G sharded)}

The same pooling at $1/1200^\circ$ ($\approx 93$\,m) over declared regions:
first the deforestation frontiers (the Amazon arc, the Congo basin, insular
Southeast Asia, $\approx 5$\,M\,km$^2$), then the $\pm 100$\,km coastal band
of family 1.0.tf. Regions: $5 \times 10^8$ cells $\times$ 65 $\times$ 1.2
$\approx 39$\,GB per year, \textbf{$\approx 0.35$\,TB} for nine years; the
coastal band (taken as 30\,\% of land, $\approx 8 \times 10^9$ cells at this
grid) would be $\approx 0.6$\,TB per year, $\approx 5.7$\,TB in all, and is
phase C. The region list is a registry field, and a
new region is a new build of that region alone.

\subsection{\code{aef10} --- Google's original 10\,m files, read in place by byte range rather than copied}\label{sec:aef10}

The 10\,m tiles are not copied. The registry carries a group of layout
\code{"external-cog"}: the mirror's base URL, the per-tile index
(\code{aef\_index.parquet}, 70\,MB, re-hosted on the Hub with its sha256),
the projection per zone, the companding formula, and the bottom-up
orientation flag. A read is a COG-internal byte range on an anonymous
bucket, the same operation the family-7 globe layer (the web globe's view of the global tensor, one byte range per frame) performs on the Hub (Hugging Face),
and the reader returns raw 10\,m vectors for a cone's dots on demand. This
is the family's honest position on 577\,TB: it is readable, it is not ours
to host, and the day the mirror goes away the three copied stores above are
what remains --- which is why they exist.

\begin{table}[H]
\centering\setlength{\tabcolsep}{3pt}\footnotesize
\begin{tabular}{@{}L{1.5cm}L{1.3cm}L{4.6cm}L{2.0cm}L{2.2cm}L{1.6cm}L{0.7cm}@{}}
\toprule
store & tier & content & footprint $(\log_2 f_p, \log_2 f_t)$ & rows or cells & stored & phase \\
\midrule
\code{aef1k} & G sharded & \emph{the embedding averaged per square kilometre, with its uniformity.} 64 pooled + 1 coherence, int8/uint8, $1/120^\circ$, one frame per year & $(-4.9, +6.2)$ & 270\,M land cells $\times$ 9 & $\approx 0.2$\,TB & A \\
\code{aef\_dots} & P & \emph{unaveraged 10\,m samples at fixed places.} 16 raw 10\,m vectors per $0.25^\circ$ land cell per year, float16 & $(-11.4, +6.2)$ & 43\,M rows & 7\,GB & A \\
\code{aef100r} & G sharded & \emph{the embedding averaged per 100\,m square, in chosen regions.} pooled 100\,m over regions of interest; the coastal band later & $(-8.2, +6.2)$ & $5 \times 10^8$ cells $\times$ 9 & $\approx 0.35$\,TB (regions); $\approx 5.7$\,TB (band) & B / C \\
\code{aef10} & external & \emph{Google's original 10\,m files, read in place.} the mirror's COGs by byte range; index and formula in the registry & $(-11.4, +6.2)$ & $1.4 \times 10^{12}$ px $\times$ 9 & 0 (577\,TB referenced) & A \\
\bottomrule
\end{tabular}
\caption{The four groups family 0.9.tf builds; the groups it inherits from family 1.0.tf are listed in \S\ref{sec:inherits}.}
\end{table}

\section{Missing values --- how gaps are recorded}

Nodata ($-128$ in the source) is NaN (not-a-number) in every store and never a zero
vector; a pooled cell with fewer than 10\,\% valid pixels is NaN with its
coherence NaN, and the valid fraction is written to the shard index so a
consumer can raise the threshold. Water beyond the shallow zone, the poles
above the product's coverage and the years before 2017 are absent, and
absent is what the token says ($n_R = 0$, no readings).

\section{Design decisions --- the choices made, and why}

\paragraph{Inherit, do not copy.} The two families differ in one thing,
the imagery, and the comparison they exist for is exact only if everything
else is the same bytes. Rebuilding the point networks, the LiDAR, the daily
fields, the targets and the statics for 0.9.tf would repeat 1.0.tf's phases A
and B (0.7\,TB and 2.3--2.7\,TB, nearly all of it outside the imagery), would
draw a second sample from archives that change between fetches, and would
make every difference between the two families' results partly a difference
between two builds. Inheriting by reference costs nothing, is the pattern
family 10.2 already uses for family 10.1's stores, and makes the two
registries differ in exactly the groups that are the experiment: ten
catalogue stores on one side, four embedding stores on the other. The price
is a dependency: an inherited store exists for 0.9.tf once 1.0.tf has
published it.

\paragraph{Pool by the mean of unit vectors, and keep its length.} The
producer says the vectors are linearly composable and that the dot product
is the similarity; the mean of unit vectors is then the natural pooled
representative, its length the natural measure of agreement. Any other
pooling (a medoid, a random pixel, a principal direction) discards either
the ``what'' or the ``how uniform''. The dots store keeps raw samples beside
the mean so the pooling choice is testable rather than assumed.

\paragraph{Strict availability by year end.} The one rule that must not
bend: an embedding of year $Y$ is not visible to an anchor inside $Y$. The
cost is a lag of up to a year on the most recent land state, which is real
and is the price of an annual product; the alternative is a leak that would
make every land read-out look better than it is.

\paragraph{Build once, next to the data.} Producing \code{aef1k} and
\code{aef\_dots} means reading all 577\,TB once. On a rented box at
2\,Gbps that is 27 days and the transfer risk of the family-10.1
\code{slatrack} night (when that 67\,GB sea-level store crawled up from a slow rented host); on a spot instance (a discounted cloud machine the provider can reclaim) in the mirror's own region
(us-west-2, Amazon's Oregon region; in-region S3 reads free, anonymous) it is about five days at
10\,Gbps and \$50--150 of compute, writing year-parts to the Hub as it goes
with \code{done.json} (the completion marker) last, the family-10 lane discipline. Pick: in-region,
one lane per year, the first year measured before the other eight are
dispatched.

\paragraph{Reference the original rather than copy it.} 577\,TB is about
seventy times the largest store any family plans and is CC BY 4.0 on two
public buckets;
copying it would cost more than every other family together and add
nothing. The referenced group is the honest shape: the registry says where
the bytes are and how to read them, and the copied stores are what the
programme guarantees.

\paragraph{Why not stop at 1\,km.} A kilometre is the family-7 resolution's
neighbour, not the point. The dots store (7\,GB) and the regional 100\,m
store are what let the programme ask whether the land targets need 10\,m
inputs, at a cost that is a rounding error next to the raw families.

\paragraph{What this family cannot do, said plainly.} It cannot see a
five-day event, cannot be extended to NISAR or to any sensor Google did not
train on, cannot go back before 2017, and cannot be re-run over a chosen
support window; if Google changes or stops the product, the copied stores
freeze at the last year. Those are the reasons family 1.0.tf exists, and
the two are compared, not merged.

\section{Validation --- the checks every store must pass}

Every store as in family 10.1 (hashes, $N$ from file sizes, bin identity,
CSR --- compressed-sparse-row --- index) and family 1.gf (shards decompressing to declared shapes); plus
three checks specific to this source, each on every tile read: the
de-quantised vectors are unit length to within $10^{-2}$ --- the int8
companding itself perturbs the length by up to $8.5 \times 10^{-3}$, so a
tighter tolerance would reject correct data, and a wrong companding or a
wrong orientation breaks it by far more (the encoder must also clip to
$\pm 127$, since $|x| = 1$ maps to 128); the bottom-up
orientation is asserted by comparing a tile's corner against the index
polygon; and the per-year tile count and byte total equal the listing in
the table in \S3, so a truncated year is refused by name. For the inherited groups one check replaces the build: every file's sha256 in 0.9.tf's registry equals the entry in 1.0.tf's, compared file by file, so an inherited store that drifted is refused by name.

\section{Limitations --- what this family cannot do}

Annual, closed, 2017 onwards, terrestrial and shallow water only, and a
learned summary whose failure modes are the producer's (transfer across
regions and years, short events averaged out). The support window of the
released layers is not published, so ``what observations went into year
$Y$'' is not knowable pixel by pixel. The regional list for \code{aef100r} is
a choice, and a region not on it has only the 1\,km field and the dots.
The inherited stores are exactly as good as 1.0.tf's builds, and appear only as 1.0.tf publishes them. Nothing is dispatched on an unmeasured size: the first year's build
measures the pooling throughput, the coherence distribution and the dots'
bytes before the rest.

\section{Next steps: the ladder, the probe and the falsifiers --- build order, the test run, and the tests that decide what the family is worth}

\paragraph{Phase A.} \code{aef1k}, \code{aef\_dots}, the \code{aef10}
reference group and the registry with its \code{inherits} block, which costs no bytes --- each inherited store is listed as 1.0.tf publishes it; $\approx 0.2$\,TB; one in-region lane per
year. \textbf{Phase B.} \code{aef100r} over the deforestation frontiers.
\textbf{Phase C.} the coastal band at 100\,m, once E-078c's coastal cells
exist to compare against.

\paragraph{The probe (a first test run through the real code).} One UTM zone of one year (about 280 tiles, 0.5\,TB)
through the real pooling code: throughput, the valid fraction, the
coherence histogram, the orientation check, and the bytes of the resulting
shards --- the numbers that replace this note's before the nine-year build.
The Hugging Face sample of Major-TOM, an open Earth-observation dataset project (\code{Major-TOM/Core-AlphaEarth-Embeddings}, 62.5\,k
grid cells, 2017--2023) is the unit-test fixture.

\paragraph{Three falsifiers (tests that can show the family is not worth its cost), three seeds (independent training runs) each.}
\begin{enumerate}\itemsep2pt
\item \textbf{Against the programme's own codec.} On the land pixels of
E-078c's $\approx 300$ coastal cells, tokens from \code{aef\_dots} and
\code{aef1k} against the local codec's cached tokens, on the \emph{annual}
land targets in the band (loss year, land cover, biomass) --- not on E-078c's
own pentad ocean targets, which the embedding does not cover and whose
cadence it cannot match. If AlphaEarth matches or beats the codec there,
the codec is not needed for land at annual cadence and E-078c's scope
narrows to what the embedding cannot see.
\item \textbf{Against the coarse tensor, on the land targets.} Next year's
forest alerts, loss year, burned area and water extent (family 1.0.tf's
target stores) predicted from family 7.2 (the global $0.25^\circ$ five-day tensor) alone versus family 7.2 plus this
family's AlphaEarth tokens, at the targets' own footprints. Parity means 10\,m
information does not reach those targets through an annual summary.
\item \textbf{What the annual cadence costs.} The same targets from the AlphaEarth
tokens plus the daily fields this family inherits from 1.0.tf (temperature, reflectance, snow,
fire) versus the AlphaEarth tokens alone. The gap is the value of the five-day rule on
land, measured rather than argued, and it is the number that decides how
much of family 1.0.tf's image tier to build.
\end{enumerate}

\section{Decisions for Chris --- open choices, each with this note's pick}

\begin{enumerate}\itemsep2pt
\item \textbf{The name}: 0.9.tf as given, the number reading as ``an
embedding in place of the raw record''. Pick: as given.
\item \textbf{Inherit by reference}: every family 1.0.tf store except the
ten image catalogues and what is derived from them is listed in 0.9.tf's
registry with its 1.0.tf path and sha256, and none is rebuilt, so 0.9.tf's
own build is the four AlphaEarth stores (\S\ref{sec:inherits}). Pick:
inherit by reference.
\item \textbf{Strict availability by year end.} Pick: strict.
\item \textbf{Pooling}: renormalised mean plus coherence. Pick: yes, with
the dots store as the check.
\item \textbf{Where to build}: a spot instance in the mirror's region
rather than a rented box. Pick: in-region.
\item \textbf{The regional list for 100\,m}: the three deforestation
frontiers first. Pick: yes.
\item \textbf{Custom Satellite Embeddings}: the academic programme for
sub-annual embeddings closes 2026-10-15; applying costs a form and would
make falsifier 3 answerable at five-day cadence on the same model. Pick:
apply; a granted sample becomes the private-track store \code{aef\_custom}.
\end{enumerate}

\section*{Terms used above, in one line each}
\footnotesize
\textbf{AlphaEarth Foundations (AEF)} Google DeepMind's model that turns Sentinel-1, Sentinel-2 and Landsat image sequences into one 64-number description of each 10\,m of land per year, trained to reproduce radar, LiDAR, climate and land-cover targets; \textbf{Satellite Embedding V1} the published output. \textbf{Embedding} a learned vector summarising what was observed. \textbf{Unit-length / $S^{63}$} vectors of length one on a 63-dimensional sphere, so their dot product is a similarity. \textbf{int8 companding} storing a value in one byte on a square-root scale. \textbf{COG / VRT / UTM} cloud-optimised GeoTIFF (readable by byte range), a GDAL virtual-raster wrapper, and the projected map zones the tiles are in. \textbf{Tier G / P} family 10's gridded and point storage tiers. \textbf{Coherence} here, the length of the mean of unit vectors in a cell: 1 if all agree, 0 if they cancel. \textbf{Sentinel-1/2, Landsat} the model's inputs: C-band radar and optical imagery from ESA and optical from NASA/USGS. \textbf{PALSAR-2, GEDI, ERA5-Land, GRACE, GLO-30, NLCD, GBIF} its training targets: L-band radar from JAXA, spaceborne LiDAR, a land climate reanalysis, satellite gravity, an elevation model, the US land-cover map, and species-occurrence records. \textbf{SST} sea-surface temperature. \textbf{NaN} not-a-number, the stored value for ``not measured''. \textbf{CSR} compressed sparse row, the bin index of a point store. \textbf{Earth Engine} Google's hosted geospatial platform, which needs an account. \textbf{S3 / us-west-2 / spot instance} Amazon's object storage, its Oregon region where the mirror lives, and a discounted pre-emptible virtual machine. \textbf{sha256} the file hash every store carries. \textbf{NISAR} the NASA--ISRO L-band radar launched 2025, not among the inputs. \textbf{E-078 / E-078c} the multi-granularity design and its first coastal tile-codec experiment. \textbf{Source Cooperative / AWS Open Data / GCS} the mirror's host, the sponsorship it runs under, and Google Cloud Storage. \textbf{Major-TOM} an open Earth-observation dataset project whose Hugging Face sample of the embeddings serves as a test fixture. \textbf{Public / private track} a store anyone may read from the public repository, or one held for this project's own use because its licence forbids redistribution (family 1.gf \S4).
\normalsize

\begin{thebibliography}{9}\small
\bibitem{aef} Brown, C.\ F.\ et al.\ (2025). AlphaEarth Foundations: an embedding field model for accurate and efficient global mapping from sparse label data. arXiv:2507.22291 (v2, 2025-09-08).
\bibitem{aefee} Google Earth Engine / Google DeepMind (2026). \code{GOOGLE/SATELLITE\_EMBEDDING/V1/ANNUAL}, years 2017--2025. CC BY 4.0.
\bibitem{aefgcs} Google (2026). AlphaEarth Foundations Satellite Embeddings on Google Cloud Storage: read-me, index files and de-quantisation. Provider-pays since July 2026.
\bibitem{aefsc} Taylor Geospatial Engine (2026). \code{tge-labs/aef}, Source Cooperative mirror of Satellite Embedding V1, 2017--2025.
\bibitem{f1tf} blauewelt/earth (2026). Family 1.0.tf: land and coast at ten kilometres and five days. Design note, 16 September 2026.
\bibitem{f10} blauewelt/earth (2026). Family 10.1: the granularity-aware tensor, as built. Design note, 16 September 2026.
\bibitem{e078} blauewelt/earth (2026). E-078: family 10 --- one schema for every granularity. \code{ml/plans/E078\_multi\_granularity.md}.
\end{thebibliography}

\end{document}
